\section{Maximality, uniform path limits, and a common envelope}\label{sec:selection}
Fix an almost-sure maximal $h\in\clprob{K_1}$. Choose $1$-admissible integral
gains $Y_n$ with their own terminal values satisfying $Y_n(\infty)\to h$
almost surely. We first obtain uniform convergence of the paths over all time;
test-uniform Cauchy estimates are proved separately in the next section.

\subsection{Excluding persistent terminal improvements}
The next lemma serves the role of \cite[Lemma A3.1]{DS} in the original proof of Lemma 4.11: persistent improvement contradicts maximality, here allowing a vanishing lower error.

\begin{lemma}[Vanishing lower error and maximality]\label{lem:persistent-improvement}
Suppose $f_n$ are measurable, $f_n\to h$ almost surely, $g_n\in K_1$, and
$(f_n-g_n)^+\to0$ almost surely. There cannot exist fixed $\eta,p>0$ and
measurable events $E_n$ such that $\Pp(E_n)\geq p$ and $g_n-f_n\geq\eta$
on $E_n$ for every $n$.
\end{lemma}
\begin{proof}
Put $e_n=(f_n-g_n)^+$ and $b_n=f_n-e_n=\min(f_n,g_n)$.
Then $b_n\leq g_n$ and $b_n\to h$ almost surely. Apply forward convex
compactness to the uncorrected claims $g_n\in K_1$. For the resulting weights
$W_n$, we have $W_ng\to g\in\clprob{K_1}$ and $W_nb\to h$ almost surely.
Thus $g\geq h$, and maximality gives $g=h$ almost surely.
The nonnegative gaps $d_n=g_n-b_n$ exceed $\eta$ on $E_n$.
Their truncations $c_n=\min(d_n,\eta)$ satisfy $\E c_n\geq\eta p$ and
\[
 0\leq W_nc\leq\min(W_nd,\eta),\qquad \E(W_nc)\geq\eta p.
\]
But $W_nd=W_ng-W_nb\to0$ almost surely, so bounded convergence forces
$\E(W_nc)\to0$, a contradiction. No membership assumption on $f_n$ is used.
\end{proof}

\subsection{Chronological pasting}
The next proposition corresponds to \cite[Lemma 4.5]{DS}: maximality and time-ordered pasting give Cauchy estimates for the all-time supremum.

\begin{proposition}\label{prop:all-time-cauchy}
After extraction, the approximating gains satisfy
\[
 \sup_{t\geq0}|Y_n(t)-Y_k(t)|\longrightarrow0
 \quad\text{in probability}\quad(n,k\longrightarrow\infty).
\]
\end{proposition}
\begin{proof}
Failure supplies pairs in every tail whose gap exceeds a fixed positive
amount with probability bounded below. Right continuity reduces the all-time
supremum to a countable time set; approximate its event by a finite subset,
order it chronologically, and choose one of the two orientations with
positive probability bounded below. Hold the higher candidate until the first
grid time detecting the positive gap, then follow the other candidate's
increments. The decision is measurable at the switching time, so \contractref{M5}
supplies this strategy and its specified terminal value. The positive gap preserves the lower bound $-1$.
The terminal gain equals the other candidate's terminal value plus the detected
gap when switching occurs. Both original terminal sequences converge to $h$;
therefore the pasted claims have vanishing lower error and persistent positive
improvement, contradicting the preceding lemma. The countable enumeration
only detects the supremum; trades always switch on the chronological grid.
\end{proof}
\begin{proposition}
A further subsequence converges almost surely, uniformly over all time, to a
zero-start strongly adapted c\`adl\`ag process $X$ with $X\geq-1$ and
$X_t\to h$ almost surely as $t\to\infty$.
\end{proposition}
\begin{proof}
Apply \contractref{P1} to obtain a subsequence with a strongly adapted c\`adl\`ag uniform limit.
The initial value and lower bound pass to the limit.
Finally,
\[
 |X_t-h|\leq\sup_u|X_u-Y_n(u)|+
 |Y_n(t)-Y_n(\infty)|+|Y_n(\infty)-h|.
\]
First choose large $n$, then large $t$, to identify the terminal limit of
this same $X$ with $h$.
\end{proof}

\subsection{One measure and separate gain decompositions}
Use \contractref{P2} for a further subsequence and a common finite nonnegative measurable envelope $q$.
Then \contractref{P3} supplies one equivalent measure $Q_*$ with $q\in L^2(Q_*)$.
Apply \contractref{D1} to each $Y_i$ to obtain the zero-start special decomposition
$Y_i=M_i+A_i$ of the same gain, with predictable locally finite-variation $A_i$.
Each gain's $L^2$ envelope feeds the decomposition construction separately;
the auxiliary source may depend on $i$. The next section estimates arbitrary
indexed classes, including the convex and normalized tail families themselves.
All selections retain the same measure, gains, path limit $X$, and terminal $h$.
